\scanpage{33}
% archive sheet number: 32
\begingroup
\makeatletter\tagsleft@true\makeatother
% loc: opening and (5), upper line
Notons les implications
\[
\left\{\begin{aligned}
&X\ \text{connexe}\quad\Longrightarrow\quad E\overset{\mathrm{def}}{=}\bigcap_{\varepsilon\in\omega}\deleted{\unsure{X|\Sigma}}\Sigma^\varepsilon\ \deleted{\unsure{-|=}}\deleted{\unsure{=|\varnothing}}\ne\varnothing\\
% loc: (5), lower line and left of brace
&\text{\unsure{\deleted{Si}|Si|si} les }\Sigma^\varepsilon\ \text{connexes, alors}\quad X\ \text{connexe}\ \Longleftrightarrow\ E\ne\varnothing.
\end{aligned}\right.\tag*{\shortstack[l]{(5)\\\unsure{,|'}}}
\]

% loc: (6), definition
Nous poserons
\[
\Sigma^{\unsure{\circ|o|0}\varepsilon}=\Sigma^\varepsilon\setminus E=\Sigma^\varepsilon\setminus\Sigma^\varepsilon\cap\Sigma^{\varepsilon'}=X\setminus\Sigma^{\varepsilon'}\qquad(\unsure{\omega|w}\unsure{\ne|=|\in}\{\varepsilon,\varepsilon'\})\tag*{(6)}
\]
% loc: (7), modified exponents and parenthesis
de sorte que l'on a la décomposition
\[
X=\Sigma^{\circ\deleted{\unsure{\varepsilon|2}}\varepsilon_{\deleted{\unsure{1|+}}\unsure{.|,}}}\cup E\cup\Sigma^{\circ\varepsilon'_{\deleted{\unsure{2|-}}}}\qquad(\omega=\{\unsure{\varepsilon'|\varepsilon}_{\deleted{\unsure{1|+}}},\varepsilon'_{\deleted{\unsure{2|-}}}\})\tag*{(7)}
\]
de $X$ en \deleted{trois} ens. mutuellement disjoints.

% loc: (8), bars, point and punctuation
Notons
\[
\overline{\Sigma^{\circ\varepsilon}}\cap E=\overline{\Sigma^{\circ\varepsilon}}\cap\Sigma^{\varepsilon'}=\overline{\Sigma^{\circ\varepsilon}}\setminus\Sigma^{\circ\varepsilon}=\dot{\Sigma}^{\circ,\varepsilon}\quad\unsure{;|\varepsilon,|\varepsilon|,}\tag*{(8)}
\]

% loc: two struck lines before 1.2
\deleted{On a le \unsure{cas}}\par
\deleted{\unsure{on|On} s'intéresse à la \unsure{condi|cond}}

% loc: 1.2 and (9)
1.2. Deux décompositions binaires
\[
\pi_1=\{\Sigma_{\deleted{\unsure{\varepsilon|2}}1}^\varepsilon\}_{\varepsilon\in\omega_1},\qquad\pi_2=\{\Sigma_2^\varepsilon\}_{\varepsilon\in\omega_2}\tag*{(9)}
\]
% loc: line before (10) and insertion
sont dites non \uline{\unsure{parallèles|paralèles}}, \inserted{si elles sont distinctes, et} s'il existe un \unsure{.|,} $\Sigma_1^+\in\pi_1$, et un $\Sigma_2^+\in\pi_2$, tels que l'on ait
\[
\Sigma_1^+\cap\Sigma_2^+=\varnothing\tag*{(10)}
\]
% loc: after (10), (11) and (11 bis)
condition qui équivaut à
\[
\Sigma_1^+\subset\Sigma_2^-\tag*{(11)}
\]
ou encore à
\[
\Sigma_2^+\subset\Sigma_1^-\tag*{(\unsure{11bis|11 bis|11'bis})}
\]
% loc: last line and (12)
où $\Sigma_1^-,\Sigma_{\unsure{2|1}}^-$ sont \unsure{définis|définies} par
\[
\pi_1=\{\Sigma_1^+,\Sigma_1^-\},\qquad\pi_2=\{\Sigma_2^+,\Sigma_{\unsure{2|1}}^-\}\tag*{(12)}
\]
\endgroup
